\subsection{Graph Theory}

\subsubsection{Core Concepts}


%\begin{comment}
\begin{figure}[ht]
    \centering
    \begin{tikzpicture}[scale=1, transform shape]
        \tikzstyle{every node} = [draw,shape=circle,color=blue]
        \node (v1) at (5,5) {$v_1$};
        \node (v2) at (3,3) {$v_4$};
        \node (v3) at (7,3) {$v_5$};
        \node (v7) at (11,3) {$v_6$};
        \node (v5) at (5,1) {$v_7$};
        \node (v6) at (7,-1) {$v_8$};
        \tikzstyle{every node} = [draw,shape=rectangle,color=teal]
       \node (v4) at (1,4) {$v_2$};
        \node (v11) at (11,5) {$v_3$}; 
       \draw[-{Latex[length=3mm, width=2mm]}, color=teal] (v4) to (v2);
       \draw[-{Latex[length=3mm, width=2mm]}, color=teal] (v11) to (v7);
       \draw[{Latex[length=3mm, width=2mm]}-{Latex[length=3mm, width=2mm]}, color=red, bend left] (v1) to (v7);
       \draw[-{Latex[length=3mm, width=2mm]}, color=blue] (v1) to (v2);
       \draw[-{Latex[length=3mm, width=2mm]}, color=blue] (v1) to (v3);
       \draw[-{Latex[length=3mm, width=2mm]}, color=blue] (v2) to (v5);
       \draw[-{Latex[length=3mm, width=2mm]}, color=blue] (v3) to (v5);
       \draw[-{Latex[length=3mm, width=2mm]}, color=blue] (v5) to (v6);
       \draw[-{Latex[length=3mm, width=2mm]}, color=blue] (v7) to (v6);
    \end{tikzpicture}
    \caption{Contextual/Conditional Acyclic Directed Mixed Graph (CADMG).}
    \label{fig:cadmg}
\end{figure}
%\end{comment}


\begin{Def}[Contextual/conditional directed mixed graphs (CDMG)]
    \label{def-cdmg}
    A \emph{contextual/conditional directed mixed graph (CDMG)} $G$ - per definition - consists of two (disjoint) sets of
    vertices/nodes:
    \begin{enumerate}[label=\roman*.)]
        \item $J$, whose elements are called input nodes,
        \item $V$, whose elements are called output nodes,
    \end{enumerate}
%    such that $J \cup V \neq \emptyset$;
    and two (disjoint) sets of edges:
    \begin{enumerate}[resume,label=\roman*.)]
        \item $E \ins (J \cup V) \times V$ the set of directed edges ,
        \item $L \ins V \times V/((v_1,v_2) \sim (v_2,v_1)) $, the set of bi-directed edges, 
        \item[]    with: $(v_1,v_2) \in L \, \implies\, v_1\neq v_2 \land (v_2,v_1) \in L$.
    \end{enumerate}
\end{Def}

\begin{Not}
    \label{not-cdmg}
    Let $G=(J,V,E,L)$ be a CDMG.
    We will write:
    \begin{enumerate}
        \item $v \in G$ to mean $v \in J \cup V$,
        \item $v_1 \tuh v_2 \in G$ to mean $(v_1,v_2) \in E$,
        \item $v_1 \hut v_2 \in G$ to mean $(v_2,v_1) \in E$,
        \item $v_1 \huh v_2 \in G$ to mean $(v_1,v_2) \in L$,
        \item $v_1 \suh v_2 \in G$ to mean that either $v_1 \tuh v_2 \in G$ or $v_1 \huh v_2 \in G$.
        \item $v_1 \hus v_2 \in G$ to mean that either $v_1 \hut v_2 \in G$ or $v_1 \huh v_2 \in G$.
        \item $v_1 \sus v_2 \in G$ to mean that either $v_1 \tuh v_2 \in G$ or $v_1 \hut v_2 \in G$ or $v_1 \huh v_2 \in G$.
        \item etc.
    \end{enumerate}
    The star on the arrow end here stands for a placeholder to mean: ``arrow head or tail''.
\end{Not}

\begin{Rem}
With the notations \ref{not-cdmg} the restrictions in definition \ref{def-cdmg} mean that the nodes $j \in J$ will not have any arrow heads pointing towards them: $ j \hus v \notin G$. Nodes $j \in J$ can only point towards nodes $v \in V$: edges $j \tuh v$ are allowed.
\end{Rem}

\begin{Def}[Walks]
     Let $G=(J,V,E,L)$ be a CDMG and $v,w \in G$.
\begin{enumerate}
    \item A \emph{walk} from $v$ to $w$ in $G$ is a finite sequence of nodes and edges 
        \[v=v_0 \sus v_1 \sus  \cdots \sus v_{n-1} \sus v_n=w\] 
        in $G$ for some $n \ge 0$, i.e.\ such that for every $k=1,\dots,n$ 
        we have that $v_{k-1} \sus v_k \in G$, and with $v_0=v$ and $v_n=w$.
    \item[] In the definition of walk the appearance of the same nodes several times is allowed. 
            Also  the \emph{trivial walk} consisting of a single node $v_0 \in G$ is allowed as well (if $v=w$).
    \item A \emph{directed walk} from $v$ to $w$ in $G$ is of the form:
        \[v=v_0 \tuh v_1 \tuh  \cdots \tuh v_{n-1} \tuh v_n=w,\]
        for some $n \ge 0$,
        where all arrow heads point in the direction of $w$ and there are no arrow heads pointing back.
   % \item[] Directed walks exclude the trivial walk per definition.
    \item A \emph{bi-directed walk} from $v$ to $w$ in $G$ is of the form:
        \[v=v_0 \huh v_1 \huh  \cdots \huh v_{n-1} \huh v_n=w,\]
        for some $n \ge 0$, where all edges are bi-directed.
    \item A \emph{collider walk} from $v$ to $w$ in $G$ is of the form:
        \[v=v_0 \suh v_1 \huh  \cdots \huh v_{n-1} \hus v_n=w,\]
        for some $n \ge 0$, where all nodes in between $v$ and $w$ have two arrow heads pointing towards them (aka collider).
        Note that for $n=1$ this reads: $v\sus w \in G$.
    \item A walk is called \emph{path} if all occurring %intermediate?
        nodes 
        are pairwise different.
    \item A \emph{trek} from $v$ to $w$ in $G$ is a path that is either trivial or of the form:
        \[v=v_0 \hut v_1 \hut  \cdots \hut v_{k-1} \hus v_k  \tuh \cdots \tuh v_{n-1} \tuh v_n=w,\]
        with all nodes $v_i \in V$ for $i=0,\dots,n$,
   %             \[v=v_0 \hut v_1 \hut  \cdots \hut v_{k-1} \huh v_k  \tuh \cdots \tuh v_{n-1} \tuh v_n=w,\]
          i.e.\ where every node has at most one arrow head pointing towards it and 
          the endnodes have an arrow head pointing towards them. 
          %(Usually treks are only considered for $v \neq w$.)
  \end{enumerate}
\end{Def}



\begin{Def}[Family relationships]
     Let $G=(J,V,E,L)$ be a CDMG, $v,w \in V$ and $A \ins J \cup V$ a subset of nodes.
     We then define:
     \begin{enumerate}
         \item The set of \emph{parents} of $v$ in $G$: \[\Pa^G(v):=\lC w \in G\,|\, w \tuh v \in G \rC.\]
         \item[] The set of \emph{parents} of $A$ in $G$: \[\Pa^G(A):= \bigcup_{v \in A} \Pa^G(v).\]
         \item The set of \emph{children} of $v$ in $G$: \[\Ch^G(v):=\lC w \in G\,|\, v \tuh w \in G \rC.\]
         \item[] The set of \emph{children} of $A$ in $G$: \[\Ch^G(A):= \bigcup_{v \in A} \Ch^G(v).\]
         \item The set of \emph{siblings} of $v$ in $G$: \[\Sib^G(v):=\lC w \in G\,|\, v \huh w \in G \rC.\]
         %\item[] The set of \emph{children} of $A$ in $G$: \[\Ch^G(A):= \bigcup_{v \in A} \Ch^G(v).\]
         \item The set of \emph{ancestors} of $v$ in $G$: 
             \[\Anc^G(v):=\lC w \in G\,|\,\exists \text{ directed walk: } w \tuh \cdots \tuh v \in G \rC.\]
             Note: $v \in \Anc^G(v)$.
         \item[] The set of \emph{ancestors} of $A$ in $G$: \[\Anc^G(A):= \bigcup_{v \in A} \Anc^G(v).\]
             Note: $A \ins \Anc^G(A)$.
         \item The set of \emph{descendants} of $v$ in $G$: 
             \[\Desc^G(v):=\lC w \in G\,|\,\exists \text{ directed walk: } v \tuh \cdots \tuh w \in G \rC.\]
             Note: $v \in \Desc^G(v)$.
         \item[] The set of \emph{descendants} of $A$ in $G$: \[\Desc^G(A):= \bigcup_{v \in A} \Desc^G(v).\]
             Note: $A \ins \Desc^G(A)$.
         \item The set of \emph{non-descendants} of $A$ in $G$:
             \[ \NonDesc^G(A) := (J \cup V) \sm \Desc^G(A).\]
         \item The \emph{strongly connected component} of $v$ in $G$:
             \[ \Sc^G(v) := \Anc^G(v) \cap \Desc^G(v).\]
             Note: $v \in \Sc^G(v)$.
         \item[] The (union of) \emph{strongly connected components} of $A$ in $G$:
             \[ \Sc^G(A) := \bigcup_{v \in A} \Sc^G(v).\]
             Note: $A \ins \Sc^G(A)$.
         \item The \emph{district} of $v$ in $G$:
             \[\Dist^G(v) := \lC w \in G\,|\, \exists \text{ bi-directed walk: } 
             v \huh v_1 \huh \cdots \huh v_{n-1} \huh w \in G \rC.\]
             Note: $v\in \Dist^G(v)$.
         \item[] The \emph{district} of $A$ in $G$:
             \[\Dist^G(A) := \bigcup_{v \in A} \Dist^G(v).\]
             Note: $A \ins \Dist^G(A)$.
         \item The d-\emph{Markov blanket} of $v$ in $G$:
             \[\MBl^G_d(v):=\lC w \in G\,|\,\exists \text{ collider walk: } 
             v \suh v_1 \huh \cdots \huh v_{n-1} \hus w \in G \rC \sm \{v\}.\]
             Note: $\lp\Pa^G(v) \cup \Ch^G(v) \cup \Pa^G(\Ch^G(v))\rp \sm \{v\} \ins \MBl^G_d(v)$.
         \item[] The d-\emph{Markov blanket} of $A$ in $G$:
             \[\MBl^G_d(A) := \bigcup_{v \in A} \MBl^G_d(v) \sm A.\]
         \item The $\sigma$-\emph{Markov blanket} of $A$ in $G$:
             \[\MBl^G_\sigma(A):= \MBl^{G^\acy}_d(A), \]
             where $G^\acy$ is an acyclification of $G$, a notion defined below in Definition~\ref{def:cdg_acyclification}.
     \end{enumerate}
\end{Def}


\begin{Def}[Acyclicity]
    \label{def-acylic}
    A  CDMG  $G=(J,V,E,L)$  is called \emph{acyclic} if there does not exist 
    any non-trivial directed walk from $v$ to itself in $G$ for any node $v \in G$.
\end{Def}

\begin{Def}
    A  Contextual/Conditional Directed Mixed Graph (CDMG)  $G=(J,V,E,L)$  is called:
   \begin{enumerate}
       \item Contextual/Conditional Acyclic Directed Mixed Graph (CADMG) if $G$ is acyclic.
       \item Directed Mixed Graph (DMG) if $J = \emptyset$.
       \item Acyclic Directed Mixed Graph (ADMG) if $G$ is acyclic and $J = \emptyset$.
       \item Contextual/Conditional Directed Graph (CDG) if $L = \emptyset$.
       \item Directed Graph (DG) if $J = \emptyset$ and $L = \emptyset$.
       \item Contextual/Conditional Directed Acyclic Graph (CDAG) if $G$ is acyclic and $L = \emptyset$.
       \item Directed Acyclic Graph (DAG) if $G$ is acyclic, $J=\emptyset$ and $L = \emptyset$.
   \end{enumerate}
\end{Def}

\begin{Def}[Topological order]
    Let $G=(J,V,E,L)$ be a CDMG.
    A \emph{topological order} of $G$ is a total order $<$ of $J \cup V$ such that for all $v,w \in G$:
    \[ v \in \Pa^G(w) \; \implies \; v < w.\]
    Equivalently, it can be described as an indexing of the nodes $J \cup V = \{v_1,\dots,v_K\}$ where parents always precede their children.
\end{Def}

\begin{Lem}
        A CDMG  $G=(J,V,E,L)$ is acyclic if and only if it has a topological order.
\end{Lem}

\begin{Def}[Predecessors]
   Let $G=(J,V,E,L)$ be a CDMG and $<$ a total order of $J \cup V$.
   The set of \emph{predecessors} of $v$ in $G$ are:
   \[\Pred^G_<(v) : = \lC w\in G\,|\, w < v \rC.\]
   We also put:
   \[\Pred^G_\le(v) : = \lC w\in G\,|\, w < v \rC \cup \{v\}.\]
\end{Def}


\subsubsection{Operations on Graphs}

%\paragraph{Interventions on Graphs}
\paragraph{Hard Interventions on Graphs}

\begin{Def}[Hard intervention on CDMGs]\label{def:G_hard_intervention}
  Let $G=(J,V,E,L)$ be a CDMG and $W \ins J \cup V$ a subset of nodes.
  The \emph{intervened CDMG} w.r.t.\ $W$ of $G$ is the CDMG: 
  $G_{\doit(W)}=(J_{\doit(W)},V_{\doit(W)},E_{\doit(W)},L_{\doit(W)})$, where:
  \begin{enumerate}[label=\roman*.)]
      \item $J_{\doit(W)}:= J \cup W$,
      \item $V_{\doit(W)}:= V \sm W$,
      \item $E_{\doit(W)}:= E \sm \lC v \tuh w \,|\, v \in G, w  \in W  \rC$,
      \item $L_{\doit(W)}:= L \sm \lC v \huh w\,|\, v \in G, w \in W \rC$,
  \end{enumerate}
  where we turn all nodes from $W$ into input nodes and remove all edges with arrow heads towards nodes from $W$.
\end{Def}

\begin{Rem}
    If $G$ is acyclic then also $G_{\doit(W)}$ is acyclic and a topological order for $G$ is also one for $G_{\doit(W)}$.
\end{Rem}

\begin{Lem}[Hard interventions commute]
      Let $G=(J,V,E,L)$ be a CDMG and $W_1, W_2 \ins J \cup V$ two  subsets of nodes from $G$. 
      Then we have:
      \[ \lp G_{\doit(W_1)} \rp_{\doit(W_2)} = \lp G_{\doit(W_2)} \rp_{\doit(W_1)} =  G_{\doit(W_1 \cup W_2)}. \]
\end{Lem}

\paragraph{Soft Interventions on Graphs}

\begin{Def}[Extending CDMGs with soft intervention nodes]
  Let $G=(J,V,E,L)$ be a CDMG and $W \ins J \cup V$ a subset of nodes.
  The \emph{extended CDMG} of $G$ w.r.t.\ nodes $W \ins J \cup V$  and corresponding \emph{soft intervention nodes} 
  $I_W = \{I_w \,|\, w \in W\}$ %(with $I_j :=j$, for $j \in J$) 
  is the CDMG: 
  $G_{\doit(I_W)}=(J_{\doit(I_W)},V_{\doit(I_W)},E_{\doit(I_W)},L_{\doit(I_W)})$, where:
  \begin{enumerate}[label=\roman*.)]
      \item $J_{\doit(I_W)}:= J\, \dot{\cup} \, \lC I_w\,|\, w \in W \rC$, %not W \sm J, makes weak and hard interventions commute
      \item $V_{\doit(I_W)}:= V$,
      \item $E_{\doit(I_W)}:= E\, \dot{\cup} \, \lC I_w \tuh w \,|\, w \in W\sm J \rC$,
      \item $L_{\doit(I_W)}:= L$,
  \end{enumerate}
  where we just add edges $I_w \tuh w$ for $ w \in W \sm J$, where $I_w$ will represent soft intervention nodes. 
\end{Def}

\begin{Rem}
    If $G$ is acyclic then also $G_{\doit(I_W)}$ is acyclic and a topological order for for $G_{\doit(I_W)}$ is also one for $G$.
    Any topological order of $G$ can be extended to one for $G_{\doit(I_W)}$ by ordering all the $I_w$ nodes first.
\end{Rem}

\begin{Lem}[Hard and soft interventions commute]
      Let $G=(J,V,E,L)$ be a CDMG and $W_1, W_2 \ins J \cup V$ two  subsets of nodes from $G$.
%      such that $W_1 \cap W_2 \ins J$.
      Then we have:
      \[ \lp G_{\doit(I_{W_1})} \rp_{\doit(I_{W_2})} = \lp G_{\doit(I_{W_2})} \rp_{\doit(I_{W_1})} 
      =  G_{\doit(I_{W_1 \cup W_2})}. \]
      We also have:
      \[ \lp G_{\doit(I_{W_1})} \rp_{\doit(W_2)} = \lp G_{\doit(W_2)} \rp_{\doit(I_{W_1})} =  G_{\doit(I_{W_1}, W_2)}. \]
\end{Lem}

\paragraph{Marginalization of Graphs}

\begin{Def}[Marginalization aka latent projection on CDMGs]\label{def:G_marginalization}
  Let $G=(J,V,E,L)$ be a CDMG and $W \ins V$ a subset of nodes from $V$.
  Then the \emph{marginalization of $G$ w.r.t.\ $W$} or the \emph{latent projection of $G$ onto $J\cup V \sm W$} is the CDMG: 
  $G^{\sm W} = (J^{\sm W},V^{\sm W},E^{\sm W},L^{\sm W})$, where::
  \begin{enumerate}[label=\roman*.)]
      \item $J^{\sm W}:= J$,
      \item $V^{\sm W}:= V \sm W$,
      \item $E^{\sm W}$ consists of all directed edges $\ul{v} \tuh \ol{v}$ with $\ul{v},\ol{v} \in J \cup V \sm W$ 
          for which there exists a directed walk in $G$:
          \[\ul{v} \tuh w_1 \tuh \cdots \tuh w_{n-1} \tuh \ol{v}, \]
          where all intermediate nodes $w_1,\dots,w_{n-1} \in W$ (if any).
      \item $L^{\sm W}$ consists of all bi-directed edges $\ul{v} \huh \ol{v}$ with 
          $\ul{v},\ol{v} \in V \sm W$, $\ul{v} \neq \ol{v}$, 
          for which there exists a trek in $G$:
          \[\ul{v} \hut w_1 \hut \cdots \hut w_{k-1} \hus w_k \tuh \cdots \tuh w_{n-1} \tuh \ol{v}, \]
          where all intermediate nodes $w_1,\dots,w_{n-1} \in W$ (if any).
      %\item $L^{\sm W}:= \lC \ul{v} \huh \ol{v}\,|\, \ul{v},\ol{v} \in V \sm W: 
      %    \exists \text{ trek: } \ul{v} \hut w_1 \cdots w_{n-1} \tuh \ol{v} \in G 
      %    \text{ with } w_1,\dots,w_{n-1} \in W \rC$,
  \end{enumerate}
\end{Def}

\begin{Rem}[Marginalization preserves ancestral relations and acyclicity]\label{rem:marg_preserves_ancestors_acyclicity}
    \begin{enumerate}
        \item For $v_1,v_2 \in G$ with $v_1,v_2 \notin W$ we have the equivalence:
            \[ v_1 \in \Anc^G(v_2)\quad \iff \quad v_1 \in \Anc^{G^{\sm W}}(v_2).   \]
        \item If the CDMG $G$ is acyclic then also $G^{\sm W}$ and a topological order of $G$ induces a 
            topological order on $G^{\sm W}$ (by just ignoring the nodes from $W$).
    \end{enumerate}
\end{Rem}

\begin{Lem}[Marginalizations commute]
      Let $G=(J,V,E,L)$ be a CDMG and $W_1, W_2 \ins V$ two disjoint subsets of nodes from $V$. 
      Then we have:
      \[ \lp G^{\sm W_1} \rp^{\sm W_2} = \lp G^{\sm W_2} \rp^{\sm W_1} =  G^{\sm (W_1 \cup W_2)}. \]
\end{Lem}

\begin{Lem}[Marginalization and intervention commute]
      Let $G=(J,V,E,L)$ be a CDMG and $W_1 \ins J \cup V$ and $W_2 \ins V$ two disjoint subsets of nodes from $G$.
      Then we have:
      \[ \lp G_{\doit(W_1)} \rp^{\sm W_2} = \lp G^{\sm W_2} \rp_{\doit(W_1)}. \]
      So we can without ambiguity also write: $G_{\doit(W_1)}^{\sm W_2}$.\\
      A similar statement holds for soft interventions.
\end{Lem}


\subsubsection{d-Separation and Sigma-Separation}

\begin{Def}[d-blocked walks]\label{def:d-blocking}
      Let $G=(J,V,E,L)$ be a CDMG and $C \ins J \cup V$ a subset of nodes and $\pi$ a walk in $G$:
      \[ \pi =\lp  v_0 \sus \cdots \sus v_n \rp.  \]
      \begin{enumerate}
          \item
              We say that the walk $\pi$ is \emph{$C$-d-blocked} or \emph{d-blocked by $C$}\footnote{The ``d'' here stands for ``directional''. d-separation was first only used for DAGs (without bi-directed edges).  For ADMGs it was then called m-separation in \cite{Ric03} to emphasize the use of the bi-directed edges. But since the notion of d-separation can be found in standard books and
              m-separation is the natural extension of d-separation we will call it d-separation again.} 
              if either:
      \begin{enumerate}[label=\roman*.)]
          \item $v_0 \in C$ or $v_n \in C$ or:
          \item there are two adjacent edges in $\pi$ of one of the following forms:
              \[\begin{array}{rccl}
           \text{ left chain: } & v_{k-1} \hut v_k \hus v_{k+1} & \text{ with } & v_k \in C,\\
           \text{ right chain: } & v_{k-1} \suh v_k \tuh v_{k+1} & \text{ with } & v_k \in C,\\
           \text{ fork: } & v_{k-1} \hut v_k \tuh v_{k+1} & \text{ with } & v_k \in C,\\
           \text{ collider: } & v_{k-1} \suh v_k \hus v_{k+1} & \text{ with } & v_k \notin C.
              \end{array}\]
      \end{enumerate}
      If we consider end nodes, left/right chain and fork as \emph{non-colliders} then we can simply say:\\
      $\pi$ is d-blocked by $C$ if it either contains a non-collider in $C$ or a collider not in $C$.
  \item We say that the walk $\pi$ is \emph{$C$-d-open} if it is not $C$-d-blocked, i.e.\ if all non-colliders are not in 
      $C$ and all collider are in $C$.
      \end{enumerate}
\end{Def}
%\Joris{
%A slightly different version is obtained by replacing the condition in the collider case.
%\begin{Def}[$d'$-blocked walks]
%  Let $G=(J,V,E,L)$ be a CDMG and $C \ins J \cup V$ a subset of nodes and $\pi$ a walk in $G$:
%  \[ \pi =\lp  v_0 \sus \cdots \sus v_n \rp.  \]
%  \begin{enumerate}
%    \item
%    We say that the walk $\pi$ is \emph{$C$-$d'$-blocked} or \emph{$d'$-blocked by $C$} if either:
%      \begin{enumerate}[label=\roman*.)]
%          \item $v_0 \in C$ or $v_n \in C$ or:
%          \item there are two adjacent edges in $\pi$ of one of the following forms:
%              \[\begin{array}{rccl}
%           \text{ left chain: } & v_{k-1} \hut v_k \hus v_{k+1} & \text{ with } & v_k \in C,\\
%           \text{ right chain: } & v_{k-1} \suh v_k \tuh v_{k+1} & \text{ with } & v_k \in C,\\
%           \text{ fork: } & v_{k-1} \hut v_k \tuh v_{k+1} & \text{ with } & v_k \in C,\\
%           \text{ collider: } & v_{k-1} \suh v_k \hus v_{k+1} & \text{ with } & v_k \notin \Anc^G(C).
%              \end{array}\]
%      \end{enumerate}
%      If we consider end nodes, left/right chain and fork as \emph{non-colliders} then we can simply say:\\
%      $\pi$ is $d'$-blocked by $C$ if it either contains a non-collider in $C$ or a collider that is not ancestor of $C$.
%  \item We say that the walk $\pi$ is \emph{$C$-$d'$-open} if it is not $C$-$d'$-blocked.
%    %, i.e.\ if all non-colliders are not in $C$ and all collider are in $C$.
%\end{enumerate}
%\end{Def}
%The following important property connects the two notions.
%\begin{Prp}
%Let $G=(J,V,E,L)$ be a CDMG, $C \ins J \cup V$ and $i, j \in V\cup J$. Then the following statements are equivalent:
%  \begin{enumerate}
%    \item there exists a $C$-$d$-open walk between $i$ and $j$ in $G$;
%    \item there exists a $C$-$d'$-open walk between $i$ and $j$ in $G$;
%    \item there exists a $C$-$d'$-open path between $i$ and $j$ in $G$.
%  \end{enumerate}
%\end{Prp}
%\begin{proof}
%TODO
%
%Since $C \subseteq \Anc^G(C)$, a collider not in $\Anc^G(C)$ is also not in $C$.
%Suppose there exists an $C$-$d$-open walk between $i$ and $j$.
%Then all non-colliders are not in $C$ and all colliders are in $C$. Hence all colliders are in $\Anc^G(C)$. Vice versa,
%suppose there exists an $C$-$d'$-open walk between $i$ and $j$. Then all non-colliders are not in $C$ and all
%colliders are in $\Anc^G(C)$. If a collider is not in $C$, we can extend the walk by replacing this collider 
%$\sto v \ots$ by the walk sequence
%$\sto v \to \dots \to c \ot \dots \ot v \ots$ with a directed path starting from $v$ and ending in the first
%node $c$ in $C$ it encounters, and from there tracing the same path backwards back to $v$. All nodes on this
%walk except for $c$ ($v$ included) are non-colliders and are not in $C$. Hence this walk is $C$-$d$-open.
%\end{proof}
%}

\begin{Def}[$\sigma$-blocked walks]\label{def:sigma_blocking}
      Let $G=(J,V,E,L)$ be a CDMG and $C \ins J \cup V$ a subset of nodes and $\pi$ a walk in $G$:
      \[ \pi =\lp  v_0 \sus \cdots \sus v_n \rp.  \]
      \begin{enumerate}
          \item
      We say that the walk $\pi$ is \emph{$C$-$\sigma$-blocked} or \emph{$\sigma$-blocked by $C$} if either:
      \begin{enumerate}[label=\roman*.)]
          \item $v_0 \in C$ or $v_n \in C$ or:
          \item there are two adjacent edges in $\pi$ of one of the following forms:
              \[\begin{array}{rccl}
  \text{ left chain: } & v_{k-1} \hut v_k \hus v_{k+1} & \text{ with } & v_k \in C \land v_k \notin \Sc^G(v_{k-1}),\\
  \text{ right chain: } & v_{k-1} \suh v_k \tuh v_{k+1} & \text{ with } & v_k \in C \land v_k \notin \Sc^G(v_{k+1}),\\
  \text{ fork: } & v_{k-1} \hut v_k \tuh v_{k+1} & \text{ with } & v_k \in C \land v_k \notin \Sc^G(v_{k-1}) \cap \Sc^G(v_{k+1}),\\
           \text{ collider: } & v_{k-1} \suh v_k \hus v_{k+1} & \text{ with } & v_k \notin C.
              \end{array}\]
      \end{enumerate}
      %If we consider end nodes, left/right chain and fork as \emph{non-colliders} then we can simply say:
      %$\pi$ is $\sigma$-blocked by $C$ if it either contains a non-collider in $C$ or a collider not in $C$.
  \item We say that the walk $\pi$ is \emph{$C$-$\sigma$-open} if it is not $C$-$\sigma$-blocked. 
   \end{enumerate}
\end{Def}
\begin{Rem}
  In words: the walk $\pi$ is \emph{$C$-$\sigma$-open} if and only if:
  \begin{enumerate}
    \item both endpoints of $\pi$ are not in $C$, and
    \item all colliders on $\pi$ are in $C$, and
    \item all non-endpoint non-colliders on $\pi$ are not in $C$ or only point to neighboring nodes in the same strongly connected component.
  \end{enumerate}
\end{Rem}

\begin{Def}[d-separation/$\sigma$-separation]
    Let $G=(J,V,E,L)$ be a CDMG and $A,B,C \ins J \cup V$ (not necessarily disjoint) subset of nodes.
    We then say that:
    \begin{enumerate}
        \item 
    \emph{$A$ is d-separated from $B$ given $C$ in $G$}, in symbols:
    \[ A \Perp^d_G B \given C, \]
    if every walk from a node in $A$ to a node in $J \cup B$ (sic!) is d-blocked by $C$. 
\item[] Otherwise we write:
    \[ A \nPerp^d_G B \given C.\]
\item       \emph{$A$ is $\sigma$-separated from $B$ given $C$ in $G$}, in symbols:
    \[ A \Perp^\sigma_G B \given C, \]
    if every walk from a node in $A$ to a node in $J \cup B$ (sic!) is $\sigma$-blocked by $C$. 
\item[] Otherwise we write:
    \[ A \nPerp^\sigma_G B \given C.\]
    \end{enumerate}
    Special case:
    \[ A \Perp_G^d B \qquad :\iff \qquad A \Perp_G^d B | \emptyset.\]
\end{Def}


\begin{Rem}
    If $G$ is acyclic then $\Sc^G(v)=\{v\}$ for all $v \in G$ and the additional conditions in ``$\sigma$-blocked'' 
    for the non-colliders are of the form $v_k \notin \Sc^G(v_{k\pm 1})=\{v_{k\pm1}\}$.
    These are then automatically satisfied, because the node $v_k$ is a parent of 
    the relevant other node $v_{k \pm1}$ and thus not equal to it. 
    So in the acyclic case d-separation and $\sigma$-separation are equivalent.
    It turns out that in the non-acyclic case $\sigma$-separation is the better concept (and as said above it also captures the acyclic case equivalently well), see \cite{FM17,FM18,FM20}.
    We will first focus on CADMGs (acyclic) for which we can restrict ourselves to the somewhat simpler d-separation.
    Later, we will pick up $\sigma$-separation again when we deal with cycles.
\end{Rem}

\begin{Lem}[d-separation/$\sigma$-separation under marginalization]\label{lem:stability_separation_marginalization}
        Let $G=(J,V,E,L)$ be a CDMG, $A,B,C \ins J \cup V$ and $D \ins V$ 
        be subset of nodes such that:
        \[ D \cap \lp A \cup B \cup C\rp = \emptyset.\]
        Then we have the equivalence:
        \[ A \Perp^d_G B \given C \qquad \iff \qquad  A \Perp^d_{G^{\sm D} }B \given C. \]
        We also have:
        \[ A \Perp^\sigma_G B \given C \qquad \iff \qquad  A \Perp^\sigma_{G^{\sm D} }B \given C. \]
\end{Lem}

%\begin{Def}[Acyclification of CDMGs]
%    \label{def-acyclification}
%Let $G=(J,V,E,L)$ be a CDMG.
%We define the \emph{acyclification} $G^\acy = (J^\acy,V^\acy,E^\acy,L^\acy)$ of $G$ as follows:
%\begin{enumerate}
%    \item  $ J^\acy:=J$ and $V^\acy :=V$.
%    \item  $v \tuh w \in E^\acy$ iff $v \notin \Sc^G(w)$ and there exists $\tilde w \in \Sc^G(w)$ such that $v \tuh \tilde w \in E$:
%        \[ E^\acy:=\lC v \tuh w \,|\, v \notin \Sc^G(w) \, \land\, \exists \tilde w \in \Sc^G(w):\; v \tuh \tilde w \in E \rC.  \]
%    \item  $v \huh w \in L^\acy$ iff $v \notin \Sc^G(w)$  and there 
%        exists $\tilde v \in \Sc^G(v)$ and $\tilde w \in \Sc^G(w)$ 
%        such that $\tilde v \huh \tilde w \in L$:
%        \[ L^\acy:=\lC v \huh w \,|\,v \notin \Sc^G(w) \, \land\, \exists \tilde v \in \Sc^G(v) \,\land\, \exists \tilde w \in \Sc^G(w):\; \tilde v \huh \tilde w \in L \rC.  \]
%\end{enumerate}
%In other words, the acyclification draws edges $ v \tuh w \in \Sc^G(w)$ from outside a strongly connected component to every node
%$\tilde w \in \Sc^G(w)$ of that strongly connected component, similarly for bi-directed edges, 
%and then removes all edges inside each strongly connected component of $G$. 
%Since strongly connected components can themselves ordered in an acyclic fashion the resulting graph $G^\acy$ is acyclic, thus a CADMG.
%\end{Def}
%
%\begin{Thm}
%    \label{thm-acy-sigma-sep}
%    Let $G=(J,V,E,L)$ be a CDMG and $G^\acy$ its acyclification.
%    For $A,B,C \ins J \cup V$ we then have the equivalence:
%    \[ A \Perp^\sigma_G B \given C  \qquad\iff\qquad A \Perp^\sigma_{G^\acy} B \given C \qquad\iff\qquad A \Perp^d_{G^\acy} B \given C.  \]
%\end{Thm}

\paragraph{Separoid Axioms for d-/Sigma-Separation}

\begin{DefThm}[(Asymmetric) separoid axioms for d-separation/$\sigma$-separation]
    \label{d-sep-separoid-axioms}
    Let $G=(J,V,E,L)$ be a CDMG and $A,B,C,D \ins J \cup V$ subset of nodes.
Then the ternary relations $\Perp = \Perp_G^d$ and $\Perp = \Perp_G^\sigma$ satisfy the following rules:
\begin{enumerate}[label=\alph*)]
    \item Extended Left Redundancy: 
    \item[] $D \ins A \implies D \Perp  B \given A$.
    \item $J$-Restricted Right Redundancy: 
    \item[] $A \Perp  \emptyset \given C \cup J$ always holds.
    \item $J$-Inverted Right Decomposition: 
    \item[] $A \Perp  B \given C \implies A \Perp  J \cup B \given C$.
	\item Left Decomposition:  
	\item[]$A \cup D \Perp  B \given C \implies D \Perp  B \given C$.
  \item Right Decomposition:  	
    \item[] $A \Perp  B \cup D\given C \implies A \Perp  D \given C$.
    \item Left Weak Union:
    \item[] $A \cup D \Perp  B  \given C \implies A \Perp  B \given D \cup C$.
  \item Right Weak Union:\item[] $A \Perp  B \cup D \given C \implies A \Perp  B \given D \cup C$.
\item Left Contraction:\item[] 
 $ (A \Perp  B \given D \cup C) \land (D \Perp  B \given C) \implies A \cup D \Perp  B \given C$.
\item Right Contraction:
\item[] $ (A \Perp  B \given D \cup C) \land (A \Perp  D \given C)  \implies A \Perp  B \cup D \given C$.
\item Right Cross Contraction:
\item[] $ (A \Perp  B \given D \cup C) \land (D \Perp  A \given C) \implies A  \Perp  B \cup D \given C$.
\item Flipped Left Cross Contraction: 
\item[] $ (A \Perp  B \given D \cup C) \land (B \Perp  D \given C) \implies B \Perp  A \cup D \given C$.
\end{enumerate}
In particular, we have the equivalences:
$$ (A \Perp  B \cup D \given C) \quad\iff\quad (A \Perp  B \given D \cup C) \quad\land\quad (A \Perp  D \given C)  ,$$
%and:
$$ (A \cup D \Perp  B \given C) \quad\iff\quad (A \Perp  B \given D \cup C) \quad\land\quad (D \Perp  B \given C).$$
%Finally, if $J = \emptyset$ then 
We also get:
\begin{enumerate}[resume,label=\alph*)]
    \item $J$-Restricted Symmetry: 
	\item[] $A \Perp  B \given C \cup J \implies B \Perp A \given C \cup J$. 
\end{enumerate}
For the special case of $J=\emptyset$ we have thus (Unrestricted) Symmetry.
\end{DefThm}


\begin{Rem}
Let the assumptions be like in \ref{d-sep-separoid-axioms}.
    We also have the following rules: 
    \begin{enumerate}[resume,label=\alph*)]
        \item[l)] Left Composition:
        \item[] $(A \Perp  B \given  C) \land (D \Perp  B \given C) \implies A \cup D \Perp  B \given C.$
        \item[m)] Right Composition:
        \item[] $(A \Perp  B \given  C) \land (A \Perp  D \given C) \implies A \Perp  B \cup  D \given C.$
         \item[n)] Left Intersection: If $A \cap D = \emptyset$ then:
        \item[] $(A \Perp  B \given D \cup  C) \land (D \Perp  B \given A \cup C) \implies A \cup D \Perp  B \given C.$       
        \item[m)] Right Intersection: If $B \cap D = \emptyset$ then:
        \item[] $(A \Perp  B \given D \cup  C) \land (A \Perp  D \given B \cup C) \implies A \Perp  B \cup D \given C.$
    \end{enumerate}
\end{Rem}
